% Task-list checkbox glyphs for the PDF render (REQ-04-01-03-05), injected via
% `include-in-header` in docs/pandoc/defaults.yaml.
%
% Pandoc renders Markdown task lists `[ ]` / `[x]` as U+2610 (BALLOT BOX, ☐) and
% U+2612 (BALLOT BOX WITH X, ☒). eisvogel's default Source Sans lacks those code
% points, and xelatex does no automatic font fallback, so the boxes come out as
% tofu. We remap the ballot code points to a font that carries them when the
% render image ships one, and fall back to a box drawn with base LaTeX otherwise.
%
% Deliberately package-free beyond fontspec (which eisvogel already loads): the
% pinned pandoc/extra image is a curated TeX Live subset and lacks newunicodechar
% and other add-ons, so a missing `\usepackage` would abort the whole render. The
% \fbox fallback is a no-tofu floor that needs nothing but the LaTeX kernel.

% Base-LaTeX fallback: an empty framed box, and a framed times for the checked
% state. \fbox, \vphantom, \kern and \times are all LaTeX-kernel — nothing to miss.
\protected\def\arqixBallotOpen{\fbox{\vphantom{X}\kern0.55em}}
\protected\def\arqixBallotDone{\fbox{\vphantom{X}$\times$}}

% Prefer a real ballot font when the image has one (fontspec is loaded by
% eisvogel). \IfFontExistsTF and \newfontfamily are fontspec, always available here.
\IfFontExistsTF{DejaVu Sans}{%
  \newfontfamily\arqixBallotFont{DejaVu Sans}%
  \protected\def\arqixBallotOpen{{\arqixBallotFont\symbol{"2610}}}%
  \protected\def\arqixBallotDone{{\arqixBallotFont\symbol{"2612}}}%
}{}%

% Route the ballot code points through the commands above by making them active.
\catcode`☐=\active \protected\def☐{\arqixBallotOpen}
\catcode`☑=\active \protected\def☑{\arqixBallotDone}
\catcode`☒=\active \protected\def☒{\arqixBallotDone}
